%%%%%%%%%%%%%%%%%%%%%%%%%%%%%%%%%%%%%%%%%%%%%%%%%%%%%%%%%%%%%%%%%%%%%%%%%%
% AI Engineer -- Feature Prototyping. Conversational AI powered by LLMs.
% Tailored: turn LLM capabilities into working prototypes & reference
% implementations (APIs / small services); conversational AI quality
% (accuracy, consistency, latency); prompt structures, I/O formats,
% retrieval-based systems, agent-style workflows; understand where models
% break (inconsistency, latency, unexpected outputs) & design around it;
% product-minded, Python, bridge model team <-> engineering.
% Font: Source Sans Pro. Accent: teal + purple. Links: clickable labels.
% Compile: pdflatex (2 passes).
%%%%%%%%%%%%%%%%%%%%%%%%%%%%%%%%%%%%%%%%%%%%%%%%%%%%%%%%%%%%%%%%%%%%%%%%%%

\documentclass[10pt,a4paper]{article}

\usepackage[T1]{fontenc}
\usepackage[utf8]{inputenc}
\usepackage[default]{sourcesanspro}   % professional, neutral sans (Source Sans Pro)
\usepackage[a4paper,top=1.0cm,bottom=0.8cm,left=1.5cm,right=1.5cm]{geometry}
\usepackage{titlesec}
\usepackage{enumitem}
\usepackage{xcolor}
\usepackage{hyperref}
\usepackage{microtype}

% ---- Colours (echo the HTML CV: teal primary, purple accent) ----
\definecolor{accent}{HTML}{117A8B}   % teal  (hsl 187,74%,32%)
\definecolor{company}{HTML}{6C2BB3}  % purple (hsl 270,70%,45%)
\definecolor{muted}{HTML}{555555}

% ---- Links: clickable label text, coloured teal (no raw URLs) ----
\hypersetup{colorlinks=true, urlcolor=accent, linkcolor=accent}

% ---- Remove paragraph indent, set sensible line spacing ----
\setlength{\parindent}{0pt}
\linespread{1.02}
\pagestyle{empty}

% ---- Section headings: uppercase, bold, teal, thin rule underneath ----
\titleformat{\section}
  {\normalfont\large\bfseries\color{accent}\scshape}
  {}{0em}{}[{\color{accent!35}\titlerule[1pt]}]
\titlespacing*{\section}{0pt}{6pt}{4pt}

% ---- Compact bullet list ----
\newlist{tight}{itemize}{1}
\setlist[tight]{leftmargin=1.3em,itemsep=2pt,topsep=2pt,parsep=0pt,label=\textcolor{accent}{\small$\bullet$}}

% ---- Entry: company / dates on line 1, role / location on line 2 ----
\newcommand{\entry}[4]{%
  {\color{company}\bfseries #1}\hfill {\color{muted}#2}\par
  \vspace{1pt}
  {\itshape #3}\hfill {\color{muted}\itshape #4}\par
  \vspace{3pt}%
}

% ---- Inline divider for the contact line ----
\newcommand{\sep}{\,\textcolor{accent!60}{\textbar}\,}

\begin{document}

%======================= HEADER =======================
\begin{center}
  {\fontsize{24}{26}\selectfont\bfseries Abhijeet Lokhande}\\[3pt]
  {\color{accent}\large AI Engineer: Conversational AI, LLM Prototyping \& Reference Implementations}\\[5pt]
  {\small
    London, United Kingdom \sep
    \href{mailto:abhijeet_lokhande@proton.me}{abhijeet\_lokhande@proton.me} \sep
    07480801382 \\[2pt]
    \href{https://www.linkedin.com/in/ablds/}{LinkedIn} \sep
    \href{https://github.com/abhijeetscode}{GitHub} \sep
    \href{https://huggingface.co/ab-ai}{Hugging Face}
  }
\end{center}
\vspace{2pt}

%======================= SUMMARY =======================
\section{Summary}
AI Engineer who turns LLM capabilities into working prototypes and clear reference implementations that engineering teams build on. Passionate about Conversational AI: I have shipped LLM-powered chat products, designed prompt structures, input and output formats, retrieval-based systems, and agent-style workflows, and improved feature quality for accuracy, consistency, and responsiveness. I understand how models behave in practice, where they break (inconsistency, latency, unexpected outputs), and how to design systems that work around those limits. Product-minded, fast at building in Python, and comfortable bridging the LLM team and software engineering.

%======================= SKILLS =======================
\section{Skills}
\begin{tight}
  \item[\textcolor{accent}{\small$\bullet$}] \textbf{Conversational AI:} LLM chat products, Multi-channel agents, Prompt structures \& I/O formats, Interaction patterns, Response quality (accuracy, consistency, latency)
  \item[\textcolor{accent}{\small$\bullet$}] \textbf{Prototyping \& Patterns:} Rapid prototyping, Retrieval-based systems (RAG), Agent-style workflows, Reference implementations (APIs / small services), Edge-case \& failure-mode analysis
  \item[\textcolor{accent}{\small$\bullet$}] \textbf{Engineering:} Python (fast, production-grade, TDD), FastAPI, REST APIs, PostgreSQL / MongoDB, Docker, Rust
  \item[\textcolor{accent}{\small$\bullet$}] \textbf{LLM \& Orchestration:} LLMs, Embeddings, Fine-tuning (LoRA, BERT), LangGraph, LangChain, LlamaIndex, Model Context Protocol (MCP), Agent evals
  \item[\textcolor{accent}{\small$\bullet$}] \textbf{ML:} PyTorch, TensorFlow, scikit-learn, NLP, Graph Neural Networks
\end{tight}

%======================= EXPERIENCE =======================
\section{Experience}
\entry{Built AI (FinTech \& AI Company)}{April 2022 -- Present}{AI/ML Engineer (Early Stage Hire)}{London, UK}
\begin{tight}
  \item Built and shipped Conversational AI products powered by LLMs: a client-facing chat agent now used by 5+ leading investment management firms across chat and live meetings, improving accuracy, consistency, and responsiveness through hands-on prototyping.
  \item Turned high-level product ideas into working prototypes and testable systems, designing prompt structures, input and output formats, retrieval-based (RAG) systems, and agent-style workflows in real product flows.
  \item Created reference implementations for engineering teams to build from: a production API server on FastAPI exposing 180+ AI-powered tools via the Model Context Protocol, plus reusable LLM services.
  \item Studied how models behave in practice, identifying edge cases, failure points, latency, and inconsistency early, then built evaluation frameworks and monitoring (groundedness, reliability, failure modes) to design around those limits.
  \item Built an autonomous agent that orchestrates multi-step reasoning to generate Discounted Cash Flow (DCF) models from raw data, cutting analyst modelling time by 70\%.
  \item Worked closely with product and engineering stakeholders to ensure prototypes could be built and scaled into production.
\end{tight}

\entry{King's College London}{March 2021 -- January 2022}{Research Associate}{London, UK}
\begin{tight}
  \item Built and evaluated deep learning models (Graph Convolution Networks, GANs) for drug discovery over a million-row dataset, achieving a 63\% improvement in the novelty score of generated structures.
\end{tight}

%======================= PROJECTS =======================
\section{Open-Source \& Projects}
\begin{tight}
  \item \textbf{Spec-to-Code Agent (agent-style workflow):} autonomous ReAct agent on LangGraph and the Anthropic API that plans, calls tools, and ships installable, tested packages from plain-language specs; 9/9 pass rate across a polyglot benchmark (Python, Go, Rust, Java, JS). \href{https://github.com/abhijeetscode/Spec-to-Code-Agent}{View on GitHub}
  \item \textbf{PII Detection LLMs:} fine-tuned GPT (LoRA) and BERT for PII detection and deployed them as reusable production models; collectively 67K+ downloads on HuggingFace. \href{https://huggingface.co/ab-ai/pii_model}{View on Hugging Face}
  \item \textbf{Transformer from Scratch in Rust:} GPT-style decoder-only transformer built from first principles using the Candle framework, giving a deep working grasp of how these models behave. \href{https://github.com/abhijeetscode/transformer-rust}{View on GitHub}
\end{tight}

%======================= EDUCATION =======================
\section{Education}
{\color{company}\bfseries King's College London}\hfill {\color{muted}London, UK}\par
{\itshape MSc in Data Science}\par\vspace{2pt}
{\color{company}\bfseries C-DAC, India}\hfill {\color{muted}Pune, India}\par
{\itshape Post Graduate Diploma in Advanced Computing (PG-DAC)}\par\vspace{2pt}
{\color{company}\bfseries University of Pune}\hfill {\color{muted}Pune, India}\par
{\itshape B.E. in Computer Science}\par

\end{document}
